\documentclass[11pt]{article}
\usepackage{amsmath,amssymb,amsthm}
\usepackage[margin=1in]{geometry}

\newtheorem{theorem}{Theorem}
\newtheorem{lemma}[theorem]{Lemma}

\title{Problem 527: Randomized Binary Search}
\author{Project Euler Solutions}
\date{}

\begin{document}
\maketitle

\section{Problem Statement}
A secret integer~$t$ is selected uniformly at random from $\{1, 2, \ldots, n\}$. We make guesses with feedback ($<$, $=$, $>$).

\textbf{Standard binary search} $B(n)$: always guess $g = \lfloor(L+H)/2\rfloor$.

\textbf{Randomized binary search} $R(n)$: guess $g$ uniformly from $\{L,\ldots,H\}$.

Given $B(6) = 2.33333333$ and $R(6) = 2.71666667$, find $R(10^{10}) - B(10^{10})$ rounded to 8~decimal places.

\section{Mathematical Foundation}

\begin{theorem}[Recurrence for $R(n)$]
With $R(1) = 1$ and for $n \ge 2$:
\[
R(n) = 1 + \frac{2}{n^2}\sum_{j=1}^{n-1} j\,R(j).
\]
Equivalently, after telescoping:
\[
n^2 R(n) = (2n-1) + (n^2-1)\,R(n-1).
\]
\end{theorem}

\begin{proof}
When the range has size~$n$, the guess $g$ is uniform on $\{1,\ldots,n\}$ (in relative coordinates). If $g$ lands at position~$k$:
\begin{itemize}
    \item With probability $1/n$, $g = t$ (done).
    \item With probability $(k-1)/n$, $t$ lies in the left sub-range of size $k-1$.
    \item With probability $(n-k)/n$, $t$ lies in the right sub-range of size $n-k$.
\end{itemize}
Averaging over $g$ (uniform) and $t$ (uniform):
\[
R(n) = 1 + \frac{1}{n^2}\sum_{k=1}^{n}\bigl[(k-1)R(k-1)+(n-k)R(n-k)\bigr] = 1 + \frac{2}{n^2}\sum_{j=1}^{n-1}j\,R(j).
\]
Multiplying by $n^2$ and subtracting the equation for $n-1$:
\[
n^2 R(n) - (n-1)^2 R(n-1) = (2n-1) + 2(n-1)R(n-1),
\]
which simplifies to $n^2 R(n) = (2n-1) + (n^2-1)R(n-1)$.
\end{proof}

\begin{theorem}[Standard Binary Search]
$B(n)$ equals the average depth in the implicit binary search tree on $n$~nodes:
\[
n\,B(n) = n + \lfloor(n-1)/2\rfloor\,B\bigl(\lfloor(n-1)/2\rfloor\bigr) + \lceil(n-1)/2\rceil\,B\bigl(\lceil(n-1)/2\rceil\bigr).
\]
\end{theorem}

\begin{proof}
The first guess is $m = \lfloor(1+n)/2\rfloor$, costing one guess. If $t < m$, we recurse on a range of size $m-1 = \lfloor(n-1)/2\rfloor$; if $t > m$, on a range of size $n - m = \lceil(n-1)/2\rceil$. Each sub-range is reached with probability proportional to its size, giving the stated recursion.
\end{proof}

\begin{lemma}[Asymptotic Expansion of $H_n$]
For large~$n$:
\[
H_n = \ln n + \gamma + \frac{1}{2n} - \frac{1}{12n^2} + \frac{1}{120n^4} - \cdots
\]
where $\gamma = 0.5772156649\ldots$ is the Euler--Mascheroni constant.
\end{lemma}

\begin{proof}
This is the standard Euler--Maclaurin expansion of the harmonic series. The error after $k$ correction terms is $O(n^{-(2k+1)})$, proved by bounding the remainder in the Euler--Maclaurin formula using Bernoulli number estimates.
\end{proof}

\section{Algorithm}

\begin{enumerate}
    \item Compute $B(10^{10})$ via the binary recursion in $O(\log n)$ time.
    \item Compute $R(10^{10})$ via the recurrence $n^2 R(n) = (2n-1)+(n^2-1)R(n-1)$ iteratively (or via the asymptotic solution with sufficient precision).
    \item Return $R(10^{10}) - B(10^{10})$ rounded to 8~decimal places.
\end{enumerate}

\section{Complexity Analysis}
\begin{itemize}
    \item \textbf{Time:} $O(\log n)$ for $B(n)$; $O(n)$ for $R(n)$ via direct iteration (or $O(1)$ with asymptotic expansion).
    \item \textbf{Space:} $O(\log n)$ for $B(n)$; $O(1)$ for $R(n)$.
\end{itemize}

\section{Answer}

\[
\boxed{11.92412011}
\]

\end{document}
